% S1 Appendix: tabelas suplementares. As tabelas sao geradas por scripts/build_paper_assets.py.
\documentclass[10pt,letterpaper]{article}
\usepackage[margin=0.8in]{geometry}
\usepackage{booktabs,longtable,pdflscape,caption,xurl}
\usepackage{amsmath}
\captionsetup{labelfont=bf,labelsep=period,justification=raggedright,singlelinecheck=off}
\renewcommand{\thetable}{\Alph{table}}
\begin{document}
\raggedright
{\large \textbf{S1 Appendix. Supplementary tables}}

\medskip
A zero-shot foundation model versus city-trained models for 12-month dengue forecasting in eight Brazilian state capitals: A rolling-origin benchmark. Fabiano Novaes Barcellos Filho.

\medskip
All tables are generated from the saved forecasts by \texttt{scripts/analyze\_results.py} and \texttt{scripts/build\_paper\_assets.py} (\url{https://github.com/fabianofilho/dengue-timeseries-skforecast}, tag \texttt{paper-v2}). sMAPE, symmetric mean absolute percentage error; MAE, mean absolute error; RMSE, root mean squared error. Evaluation: expanding-window rolling origin, 121 origins per city (forecast windows starting January 2014 to January 2024), 12-month horizon.

\begin{landscape}
% gerado por scripts/build_paper_assets.py -- nao editar a mao

\begin{table}[!ht]
\centering
\footnotesize
\setlength{\tabcolsep}{3pt}
\caption{\textbf{sMAPE (\%) with 95\% moving block bootstrap intervals, primary model set.}}
\begin{tabular}{lrrrrrrrr}
\toprule
City & TimesFM & CatBoost & XGBoost & RF & LightGBM & SARIMA & Prophet & Naive \\
\midrule
S\~ao Paulo & 78.0 (61.1--87.3) & \textbf{75.6 (58.4--87.3)} & 81.2 (68.3--90.5) & 77.0 (60.9--87.0) & 88.9 (81.8--93.8) & 87.8 (74.3--100.2) & 77.3 (60.0--87.8) & 77.6 (63.8--85.1) \\
Rio de Janeiro & 96.8 (84.8--109.6) & 97.5 (77.6--113.7) & 101.5 (89.0--112.8) & \textbf{93.5 (75.9--105.9)} & 101.2 (84.8--112.3) & 101.7 (87.8--111.9) & 109.6 (85.2--131.1) & 104.1 (89.3--119.0) \\
Belo Horizonte & 97.1 (85.7--111.0) & \textbf{62.0 (45.4--75.8)} & 76.4 (58.0--90.8) & 70.3 (52.0--85.8) & 92.1 (81.5--102.5) & 95.8 (87.8--107.7) & 80.7 (68.4--90.9) & 79.6 (69.6--89.5) \\
Bras\'ilia & 74.3 (62.8--86.4) & 76.6 (67.1--84.0) & 90.3 (76.5--101.1) & 81.5 (68.8--90.9) & 84.2 (74.5--90.1) & 80.2 (70.7--89.3) & 79.3 (63.8--92.8) & \textbf{74.2 (64.5--83.9)} \\
Fortaleza & \textbf{64.2 (56.6--75.0)} & 69.4 (58.6--86.6) & 79.7 (69.1--93.5) & 68.4 (56.8--84.0) & 73.7 (66.5--85.7) & 73.4 (65.3--83.8) & 69.7 (58.9--86.6) & 69.6 (59.5--86.8) \\
Recife & 67.3 (56.7--78.1) & \textbf{65.9 (50.3--80.5)} & 73.4 (58.7--86.7) & 71.0 (57.4--82.7) & 77.5 (65.6--88.9) & 87.1 (75.1--96.6) & 81.5 (61.5--103.4) & 83.3 (67.6--101.0) \\
Manaus & 56.0 (47.3--61.9) & 55.7 (47.6--61.6) & 54.8 (44.3--61.6) & 53.9 (45.5--59.4) & 62.8 (54.7--69.3) & 62.4 (50.3--71.5) & \textbf{53.8 (41.3--59.4)} & 59.1 (46.7--68.4) \\
Salvador & \textbf{68.9 (59.5--82.3)} & 71.8 (56.7--86.5) & 72.2 (57.6--84.0) & 71.0 (57.2--84.9) & 72.3 (63.6--82.1) & 79.5 (62.6--97.4) & 88.9 (68.4--108.0) & 87.2 (64.7--112.3) \\
\midrule
Mean & 75.3 & 71.8 & 78.7 & 73.3 & 81.6 & 83.5 & 80.1 & 79.3 \\
\bottomrule
\end{tabular}
\begin{flushleft} Trained models were fitted on log1p-transformed counts; TimesFM was used zero-shot on raw counts. RF, Random Forest; Naive, seasonal naive. Bold, lowest value in the row.
\end{flushleft}
\label{tabA}
\end{table}

\end{landscape}

% gerado por scripts/build_paper_assets.py -- nao editar a mao

\begin{table}[!ht]
\centering
\small
\caption{\textbf{MAE (cases per month), primary model set.}}
\begin{tabular}{lrrrrrrrr}
\toprule
City & TimesFM & CatBoost & XGBoost & RF & LightGBM & SARIMA & Prophet & Naive \\
\midrule
S\~ao Paulo & 9{,}617 & \textbf{9{,}611} & 9{,}715 & 9{,}629 & 9{,}956 & 11{,}888 & 9{,}797 & 9{,}938 \\
Rio de Janeiro & \textbf{1{,}966} & 2{,}366 & 3{,}349 & 2{,}282 & 2{,}753 & 3{,}958 & 2{,}071 & 2{,}316 \\
Belo Horizonte & 6{,}384 & \textbf{5{,}707} & 6{,}578 & 6{,}193 & 7{,}499 & 16{,}192 & 6{,}431 & 8{,}294 \\
Bras\'ilia & \textbf{3{,}823} & 3{,}965 & 4{,}290 & 4{,}090 & 4{,}095 & 5{,}558 & 4{,}023 & 4{,}319 \\
Fortaleza & \textbf{1{,}434} & 1{,}596 & 2{,}093 & 1{,}593 & 1{,}825 & 2{,}273 & 1{,}624 & 1{,}698 \\
Recife & \textbf{737} & 841 & 919 & 927 & 979 & 2{,}201 & 1{,}066 & 1{,}141 \\
Manaus & 269 & 270 & 301 & \textbf{257} & 299 & 301 & 261 & 298 \\
Salvador & \textbf{423} & 453 & 460 & 440 & 451 & 834 & 559 & 601 \\
\midrule
Mean & 3{,}082 & 3{,}101 & 3{,}463 & 3{,}176 & 3{,}482 & 5{,}401 & 3{,}229 & 3{,}576 \\
\bottomrule
\end{tabular}
\begin{flushleft} Trained models were fitted on log1p-transformed counts; TimesFM was used zero-shot on raw counts. RF, Random Forest; Naive, seasonal naive. Bold, lowest value in the row.
\end{flushleft}
\label{tabB}
\end{table}


% gerado por scripts/build_paper_assets.py -- nao editar a mao

\begin{table}[!ht]
\centering
\small
\caption{\textbf{RMSE (cases per month), primary model set.}}
\begin{tabular}{lrrrrrrrr}
\toprule
City & TimesFM & CatBoost & XGBoost & RF & LightGBM & SARIMA & Prophet & Naive \\
\midrule
S\~ao Paulo & 44{,}024 & 44{,}152 & 44{,}232 & 44{,}136 & 43{,}519 & 43{,}563 & 43{,}486 & \textbf{43{,}114} \\
Rio de Janeiro & \textbf{5{,}958} & 6{,}310 & 8{,}380 & 6{,}395 & 6{,}022 & 14{,}873 & 6{,}310 & 6{,}285 \\
Belo Horizonte & 17{,}388 & 17{,}620 & 17{,}786 & 17{,}586 & \textbf{17{,}015} & 91{,}580 & 17{,}548 & 20{,}006 \\
Bras\'ilia & 12{,}669 & 12{,}786 & 13{,}165 & 12{,}973 & 12{,}755 & 14{,}695 & \textbf{11{,}658} & 12{,}870 \\
Fortaleza & \textbf{2{,}693} & 2{,}841 & 3{,}532 & 2{,}865 & 3{,}066 & 5{,}451 & 2{,}777 & 3{,}083 \\
Recife & \textbf{1{,}474} & 1{,}616 & 1{,}731 & 1{,}666 & 1{,}685 & 10{,}907 & 1{,}984 & 2{,}005 \\
Manaus & 526 & 520 & 633 & \textbf{508} & 519 & 561 & 526 & 570 \\
Salvador & \textbf{713} & 772 & 778 & 729 & 730 & 2{,}039 & 885 & 930 \\
\midrule
Mean & 10{,}681 & 10{,}827 & 11{,}280 & 10{,}857 & 10{,}664 & 22{,}959 & 10{,}647 & 11{,}108 \\
\bottomrule
\end{tabular}
\begin{flushleft} Trained models were fitted on log1p-transformed counts; TimesFM was used zero-shot on raw counts. RF, Random Forest; Naive, seasonal naive. Bold, lowest value in the row.
\end{flushleft}
\label{tabC}
\end{table}


% gerado por scripts/build_paper_assets.py -- nao editar a mao

\begin{table}[!ht]
\centering
\small
\caption{\textbf{sMAPE (\%) with all trained models fitted on untransformed counts (sensitivity analysis).}}
\begin{tabular}{lrrrrrrrr}
\toprule
City & TimesFM & CatBoost & XGBoost & RF & LightGBM & SARIMA & Prophet & Naive \\
\midrule
S\~ao Paulo & 78.0 & 78.7 & 84.8 & 87.3 & 119.8 & 87.1 & 89.7 & \textbf{77.6} \\
Rio de Janeiro & \textbf{96.8} & 126.4 & 119.7 & 133.3 & 154.7 & 135.6 & 171.1 & 104.1 \\
Belo Horizonte & 97.1 & 95.2 & 88.7 & 90.2 & 151.4 & 134.0 & 129.0 & \textbf{79.6} \\
Bras\'ilia & 74.3 & 76.6 & 83.7 & 82.4 & 100.5 & 103.2 & 78.2 & \textbf{74.2} \\
Fortaleza & \textbf{64.2} & 77.5 & 89.4 & 98.4 & 102.1 & 87.6 & 79.2 & 69.6 \\
Recife & \textbf{67.3} & 83.3 & 83.3 & 81.2 & 111.2 & 106.1 & 102.7 & 83.3 \\
Manaus & \textbf{56.0} & 70.9 & 65.5 & 67.9 & 94.8 & 99.7 & 173.6 & 59.1 \\
Salvador & \textbf{68.9} & 76.0 & 72.1 & 73.7 & 82.2 & 91.0 & 99.5 & 87.2 \\
\midrule
Mean & 75.3 & 85.6 & 85.9 & 89.3 & 114.6 & 105.5 & 115.4 & 79.3 \\
\bottomrule
\end{tabular}
\begin{flushleft} Same design as the primary analysis, without the log1p transformation. Bold, lowest value in the row.
\end{flushleft}
\label{tabD}
\end{table}


% gerado por scripts/build_paper_assets.py -- nao editar a mao

\begin{table}[!ht]
\centering
\small
\caption{\textbf{sMAPE (\%) for the 109 origins whose forecast windows ended by December 2023 (sensitivity analysis).}}
\begin{tabular}{lrrrrrrrr}
\toprule
City & TimesFM & CatBoost & XGBoost & RF & LightGBM & SARIMA & Prophet & Naive \\
\midrule
S\~ao Paulo & 71.7 & \textbf{69.4} & 75.5 & 71.5 & 87.9 & 86.3 & 70.6 & 72.7 \\
Rio de Janeiro & 96.0 & 95.9 & 101.3 & \textbf{91.2} & 101.3 & 102.8 & 103.0 & 102.8 \\
Belo Horizonte & 95.5 & \textbf{57.1} & 72.7 & 66.5 & 89.7 & 97.1 & 76.0 & 77.7 \\
Bras\'ilia & \textbf{72.4} & 75.4 & 88.8 & 80.2 & 82.7 & 78.9 & 79.3 & 72.6 \\
Fortaleza & \textbf{66.5} & 71.9 & 79.7 & 68.3 & 74.7 & 74.2 & 71.8 & 71.9 \\
Recife & 66.7 & \textbf{65.3} & 72.4 & 69.8 & 77.0 & 86.6 & 82.7 & 84.0 \\
Manaus & 54.1 & 54.7 & 53.0 & 51.9 & 62.0 & 59.8 & \textbf{49.8} & 56.6 \\
Salvador & \textbf{69.5} & 70.8 & 71.3 & 72.1 & 72.1 & 78.8 & 84.9 & 87.4 \\
\midrule
Mean & 74.0 & 70.1 & 76.8 & 71.5 & 80.9 & 83.1 & 77.3 & 78.2 \\
\bottomrule
\end{tabular}
\begin{flushleft} Trained models were fitted on log1p-transformed counts; TimesFM was used zero-shot on raw counts. RF, Random Forest; Naive, seasonal naive. Bold, lowest value in the row.
\end{flushleft}
\label{tabE}
\end{table}


\clearpage
% gerado por scripts/build_paper_assets.py -- nao editar a mao

\begin{longtable}{llrrr}
\caption{\textbf{TimesFM versus each comparator: paired sMAPE difference and Diebold-Mariano test.} Difference = TimesFM minus comparator sMAPE, in percentage points (negative favors TimesFM), with 95\% paired moving block bootstrap interval over forecast origins (blocks of 12, 2,000 replicates). DM, Diebold-Mariano statistic on origin-level sMAPE losses with Newey-West variance (12 lags) and Harvey-Leybourne-Newbold correction; $p$ adjusted with Holm across all 56 comparisons. RF, Random Forest; Naive, seasonal naive.}\label{tabF}\\
\toprule
City & Comparator & Difference (95\% CI) & DM & $p$ (Holm) \\
\midrule
\endfirsthead
\toprule
City & Comparator & Difference (95\% CI) & DM & $p$ (Holm) \\
\midrule
\endhead
S\~ao Paulo & CatBoost & +2.4 ($-$10.2 to 17.5) & 0.31 & 1.000 \\
S\~ao Paulo & XGBoost & $-$3.2 ($-$14.3 to 8.3) & $-$0.50 & 1.000 \\
S\~ao Paulo & RF & +1.0 ($-$11.3 to 15.7) & 0.13 & 1.000 \\
S\~ao Paulo & LightGBM & $-$10.9 ($-$28.5 to 0.6) & $-$1.25 & 1.000 \\
S\~ao Paulo & SARIMA & $-$9.8 ($-$28.6 to 3.0) & $-$1.05 & 1.000 \\
S\~ao Paulo & Prophet & +0.7 ($-$5.1 to 6.7) & 0.22 & 1.000 \\
S\~ao Paulo & Naive & +0.4 ($-$14.9 to 13.6) & 0.05 & 1.000 \\
Rio de Janeiro & CatBoost & $-$0.7 ($-$14.4 to 15.2) & $-$0.08 & 1.000 \\
Rio de Janeiro & XGBoost & $-$4.7 ($-$17.2 to 7.4) & $-$0.68 & 1.000 \\
Rio de Janeiro & RF & +3.3 ($-$7.4 to 16.5) & 0.49 & 1.000 \\
Rio de Janeiro & LightGBM & $-$4.4 ($-$13.4 to 7.5) & $-$0.72 & 1.000 \\
Rio de Janeiro & SARIMA & $-$4.9 ($-$17.5 to 12.4) & $-$0.60 & 1.000 \\
Rio de Janeiro & Prophet & $-$12.8 ($-$31.6 to 8.8) & $-$1.09 & 1.000 \\
Rio de Janeiro & Naive & $-$7.3 ($-$21.0 to 5.1) & $-$1.02 & 1.000 \\
Belo Horizonte & CatBoost & +35.1 (19.9 to 56.9) & 3.38 & 0.055 \\
Belo Horizonte & XGBoost & +20.6 (6.4 to 42.9) & 2.01 & 1.000 \\
Belo Horizonte & RF & +26.8 (11.6 to 48.4) & 2.59 & 0.565 \\
Belo Horizonte & LightGBM & +4.9 ($-$11.3 to 26.5) & 0.49 & 1.000 \\
Belo Horizonte & SARIMA & +1.3 ($-$16.3 to 18.5) & 0.13 & 1.000 \\
Belo Horizonte & Prophet & +16.3 (4.2 to 35.4) & 1.87 & 1.000 \\
Belo Horizonte & Naive & +17.5 (2.6 to 36.0) & 1.86 & 1.000 \\
Bras\'ilia & CatBoost & $-$2.4 ($-$10.9 to 8.7) & $-$0.43 & 1.000 \\
Bras\'ilia & XGBoost & $-$16.0 ($-$28.7 to $-$3.0) & $-$2.29 & 1.000 \\
Bras\'ilia & RF & $-$7.3 ($-$16.3 to 3.3) & $-$1.29 & 1.000 \\
Bras\'ilia & LightGBM & $-$9.9 ($-$20.2 to 4.4) & $-$1.46 & 1.000 \\
Bras\'ilia & SARIMA & $-$6.0 ($-$19.7 to 9.0) & $-$0.76 & 1.000 \\
Bras\'ilia & Prophet & $-$5.0 ($-$17.0 to 13.2) & $-$0.56 & 1.000 \\
Bras\'ilia & Naive & +0.1 ($-$8.5 to 10.6) & 0.02 & 1.000 \\
Fortaleza & CatBoost & $-$5.2 ($-$16.5 to 5.1) & $-$0.89 & 1.000 \\
Fortaleza & XGBoost & $-$15.5 ($-$28.6 to $-$3.9) & $-$2.29 & 1.000 \\
Fortaleza & RF & $-$4.2 ($-$16.4 to 7.2) & $-$0.65 & 1.000 \\
Fortaleza & LightGBM & $-$9.5 ($-$21.2 to 1.9) & $-$1.55 & 1.000 \\
Fortaleza & SARIMA & $-$9.2 ($-$19.7 to 1.4) & $-$1.63 & 1.000 \\
Fortaleza & Prophet & $-$5.5 ($-$16.7 to 5.1) & $-$0.90 & 1.000 \\
Fortaleza & Naive & $-$5.4 ($-$18.0 to 5.7) & $-$0.84 & 1.000 \\
Recife & CatBoost & +1.4 ($-$8.7 to 13.9) & 0.22 & 1.000 \\
Recife & XGBoost & $-$6.1 ($-$13.4 to 3.1) & $-$1.33 & 1.000 \\
Recife & RF & $-$3.8 ($-$11.9 to 5.6) & $-$0.82 & 1.000 \\
Recife & LightGBM & $-$10.2 ($-$20.4 to $-$0.4) & $-$1.88 & 1.000 \\
Recife & SARIMA & $-$19.9 ($-$26.9 to $-$11.2) & $-$4.30 & 0.002 \\
Recife & Prophet & $-$14.2 ($-$30.3 to $-$0.3) & $-$1.70 & 1.000 \\
Recife & Naive & $-$16.1 ($-$28.5 to $-$6.8) & $-$2.71 & 0.404 \\
Manaus & CatBoost & +0.3 ($-$4.9 to 4.6) & 0.12 & 1.000 \\
Manaus & XGBoost & +1.3 ($-$6.1 to 7.9) & 0.33 & 1.000 \\
Manaus & RF & +2.2 ($-$2.5 to 6.9) & 0.84 & 1.000 \\
Manaus & LightGBM & $-$6.8 ($-$17.8 to 2.4) & $-$1.20 & 1.000 \\
Manaus & SARIMA & $-$6.3 ($-$15.5 to 1.3) & $-$1.37 & 1.000 \\
Manaus & Prophet & +2.2 ($-$2.4 to 10.6) & 0.57 & 1.000 \\
Manaus & Naive & $-$3.0 ($-$14.5 to 6.5) & $-$0.51 & 1.000 \\
Salvador & CatBoost & $-$3.0 ($-$11.7 to 9.3) & $-$0.50 & 1.000 \\
Salvador & XGBoost & $-$3.3 ($-$10.3 to 8.2) & $-$0.57 & 1.000 \\
Salvador & RF & $-$2.1 ($-$9.9 to 7.4) & $-$0.41 & 1.000 \\
Salvador & LightGBM & $-$3.5 ($-$7.8 to 5.0) & $-$0.92 & 1.000 \\
Salvador & SARIMA & $-$10.6 ($-$21.0 to 2.0) & $-$1.74 & 1.000 \\
Salvador & Prophet & $-$20.1 ($-$31.1 to $-$5.8) & $-$2.73 & 0.393 \\
Salvador & Naive & $-$18.3 ($-$32.6 to $-$3.2) & $-$2.21 & 1.000 \\
\bottomrule
\end{longtable}


\end{document}
